\subsection*{一步一步建立直观}

\paragraph{1. 正反馈在资源充足时近似指数}

把角速度设为零，只保留径向增长。在固定方向上，状态的大小满足“变化率与当前大小成正比”，于是呈指数增长。这可用来理解传播初期的简化机制：参与传播的人越多，单位时间新增的潜在来源也越多。图中数值没有真实人群的单位。



\paragraph{2. 区分数学相似与完整模型}

第二种情形同时带有旋转；它说明相同的指数幅度可以配上完全不同的其他机制，却不代表流行病会按这条平面螺线演化。人口流动、接触结构与易感者消耗都需要额外状态变量。本窗口仅帮助辨认线性增长因子，不能充当预测工具。



\paragraph{3. 饱和需要改变方程，而非缩放画面}

负的增长率使幅度下降，但这仍是线性指数模型。逻辑斯蒂增长则在方程中加入随存量降低的有效增长率，因而能趋近承载量。把坐标轴缩放得更大不会产生饱和；要解释增长变慢，必须指出模型中哪一项反馈发生了变化。
